\chapter{Conclusiones y Perspectivas}\label{cap.5}
\markboth{Conclusiones y Perspectivas}{Conclusiones y Perspectivas}
%
% Como hemos podido apreciar \\
%
% Este estudio nos permite tener las siguientes perspectivas:\\
En esta investigación se abordó el desarrollo de ecuaciones cuántico-relativistas para partículas escalares y fermiónicas, iniciando con la ecuación de Klein-Gordon, cuya solución reveló la presencia de densidades de probabilidad y energías negativas. Esta característica limitó su utilidad para describir sistemas físicos realistas, debido a la falta de una interpretación probabilística consistente para dichos estados. Posteriormente, se estudió la ecuación de Dirac, la cual resolvió el problema de las densidades negativas al incorporar el espín mediante matrices de Dirac, aunque mantuvo la aparición de soluciones con energías negativas. Este hecho motivó históricamente la introducción del concepto de mar de Dirac y, más adelante, la interpretación de las antipartículas dentro del marco de la teoría cuántica de campos.

Se ha demostrado que la ecuación (\ref{eq:ecuacion_de_primer_orden}) corresponde al límite no relativista de la ecuación de Dirac. Además, se verificó que a partir de dicha ecuación puede derivarse la ecuación de Schrödinger libre en tres dimensiones. Para este propósito, se introdujeron matrices simétricas nilpotentes que satisfacen la relación de anticonmutación establecida en la ecuación (\ref{eq:relacion_de_anticonmutacion}). Se construyeron explícitamente las soluciones de esta ecuación en una dimensión y se aplicaron al estudio del problema de dispersión de un electrón con espín hacia arriba frente a un potencial escalón. Se demostró que los coeficientes de transmisión y reflexión predichos por la mecánica cuántica se obtienen como la suma de los correspondientes a electrones con espín hacia arriba y hacia abajo. A partir del análisis de las gráficas, se observó que en el caso \(E > V_0\), para energías cercanas al potencial, no hay transmisión, por lo que el electrón se refleja completamente conservando su espín. Sin embargo, a medida que la energía aumenta, la probabilidad de transmisión sin flip de espín crece exponencialmente, mientras que la probabilidad de reflexión sin flip disminuye en igual proporción. Para energías comparables con la energía de reposo, se evidenció una ligera dominancia en la reflexión con inversión de espín. Para \((E<V_0\)) las particulas no se trasnmiten sino solo se reflejan.

Asimismo, se analizó el problema de una barrera de potencial finita, y se demostró que, en este caso también, la suma de los coeficientes de transmisión y reflexión para electrones con espín hacia arriba y hacia abajo reproduce los resultados establecidos por la mecánica cuántica. El análisis predice que un electrón con espín hacia arriba, al incidir sobre una barrera de potencial finita, es siempre transmitido conservando su espín. Para energías bajas, el electrón reflejado tiende a mantener su espín original. Además, para el caso \(V_0 = 1\,\mathrm{eV}\) y \(L = 1\,\mathrm{nm}\), se observó el efecto túnel cuántico, caracterizado por la disminución exponencial de la transmisión al reducirse la energía incidente. Los resultados obtenidos a partir de la ecuación (\ref{eq:version_unidimencional}) concuerdan satisfactoriamente con las expresiones bien establecidas de la mecánica cuántica. Los análisis realizados en este trabajo permiten predecir con precisión los coeficientes de reflexión y transmisión para electrones con espín hacia arriba y hacia abajo, tanto en el caso del potencial escalón como en el de la barrera finita. Se requieren pruebas experimentales que permitan contrastar las predicciones aquí obtenidas. Con base en estas conclusiones, se sugiere que la ecuación (\ref{eq:version_unidimencional}) proporciona una descripción más adecuada para partículas de espín \(1/2\) en el régimen no relativista.
